\section{Native \(\Vfour\) voltage and holonomy}

Use the blind generators as a coordinate basis:
\[
  65=(1,0),\qquad
  124=(0,1),\qquad
  326=(1,1).
\]
Choose the minimum-labeled state in each \(\Vfour\) fiber as its
origin.  Every native edge then determines the coordinate difference
between its endpoints.  The four lifts of each \(\Gfifteen\) edge give
the same difference, so a well-defined native voltage
\[
  \alpha:E(\Gfifteen)\longrightarrow\Ftwo^2
\]
is obtained.

In the registered section the thirty voltages have profile
\[
\begin{array}{c|rrrr}
\alpha(e)&(0,0)&(0,1)&(1,0)&(1,1)\\
\hline
\text{edge count}&6&16&4&4.
\end{array}
\]
The corresponding regular lift reconstructs all sixty vertices and all
\(120\) edges of \(\Gsixty\) exactly.

The cycle rank of \(\Gfifteen\) is
\[
  30-15+1=16.
\]
After tree normalization, the sixteen fundamental circuit receipts
have profile
\[
\begin{array}{c|rrrr}
h&(0,0)&(0,1)&(1,0)&(1,1)\\
\hline
\text{circuit count}&6&4&4&2.
\end{array}
\]
These receipts span all of \(\Vfour\), so the lift is connected.
Forty-five single-vertex gauge tests, covering each vertex and each
nonzero gauge value, leave the normalized circuit data unchanged.

The independently preserved certificate 033 contains a historical
thirty-edge \(\Vfour\) voltage.  We compared the derived voltage under
all \(120\) graph automorphisms of \(\Gfifteen\), all six bases in
\(\operatorname{GL}(2,2)\), and the vertex gauge determined along a
spanning tree.  Among \(720\) graph-basis comparisons, \(120\) match.
Most strongly, the identity graph labeling has exactly one match:
\[
  \text{graph map}=1,\qquad
  \text{basis map}=I_2,\qquad
  g(v)=(0,0)\quad\text{for all }v.
\]

\begin{theorem}[Native voltage recovery]
The voltage derived from the preregistered normal \(\Vfour\) action
equals certificate 033 row for row in the registered chart.  Its
regular lift is exactly \(\Gsixty\), and its holonomy image is
\(\Vfour\).
\end{theorem}
